%  Appendix A: prompt, schema and API configuration.
%  A is transcribed from the code, not typed.  The prompt is MechanicsPromptBuilder.build_regime_prompt with the 30 data lines
%  replaced by the placeholder {gex_data_table}; the schema and the request are common.REGIME_JSON_SCHEMA and common.chat_body.
%  The three typographic symbols of the prompt (U+2264, U+2265, U+2192) are written in ASCII (<=, >=, ->) in this file and
%  typeset as symbols by the "prompt" listing style, which keeps the sources ASCII.
%  Prompt check (2026-10-31): the listing of A.2, with the data lines replaced by the placeholder and the three symbols in ASCII,
%  equals common.build_prompt on 35,365 of 35,365 obfuscated gpt-6-luna requests stored in data/backups/gex_llm/*/input_*.jsonl.
%=================================================================
\section{Prompt, Schema and API Configuration}
\label{app:prompt}

This appendix reproduces the complete request sent to the model for each 30-day window, so that the experiment can be repeated from the manuscript alone. The prompt text is the output of \ident{MechanicsPromptBuilder.build_regime_prompt} (file \path{src/llm/mechanics_prompt_builder.py}); the schema and the request body are the constant \ident{REGIME_JSON_SCHEMA} and the function \ident{chat_body} of \path{scripts/validation/electronics/common.py}. The listings were generated from these objects and compared with them.

\subsection{Model and API Configuration}
\label{app:prompt:config}

% src: scripts/validation/electronics/common.py BATCH_MAX_REQUESTS=1500, MAX_COMPLETION_TOKENS=4000, RETRY_MAX_COMPLETION_TOKENS=8000; cost_ledger.yaml pricing_usd_per_mtok gpt-6-luna batch input 0.05 output 0.25 (standard 0.1/0.5); 02_decisions_log.md line 60 (model created 2026-09-14T17:15:02Z); reasoning effort medium
\begin{itemize}[leftmargin=*,labelsep=5.5mm]
\item \textbf{Model:} \ident{gpt-6-luna}, an OpenAI model alias without a dated snapshot. The creation timestamp of the model object (2026-09-14) is stored with every run and the response identifier with every answer. The published run used \ident{o4-mini-2025-04-16}.
\item \textbf{Access:} OpenAI Batch API, endpoint \path{/v1/chat/completions}, completion window of 24 hours, at most 1{,}500 requests per batch job. Each input line is one request that carries the key \texttt{<arm>|<symbol>|<window\_end>} in \ident{custom_id}; a test-retest duplicate adds the suffix \texttt{\#rep1}, and its arm is stored as \texttt{<arm>\#rep1}.
\item \textbf{Reasoning effort:} \ident{medium}, the default effort of o4-mini in the published run.
\item \textbf{Maximum completion tokens:} 4{,}000; reasoning tokens count against the limit. An answer that ends with \ident{finish_reason} = \ident{length} is requested once more with 8{,}000.
\item \textbf{Response format:} the strict JSON schema of Appendix~\ref{app:prompt:schema}. The earlier runs requested the same nine fields in the prompt text only; this is the one protocol change of the study.
\item \textbf{Not sent:} temperature, seed, system message, tools and stop sequences.
\item \textbf{Parsing:} an answer is parsed as JSON; a second attempt tolerates raw control characters inside strings, and a last fallback extracts the JSON object from free text. The parse status of every answer (\ident{ok}, \ident{fallback}, \ident{failed} or \ident{missing}) is stored.
\item \textbf{Prices:} USD 0.05 per million input tokens and USD 0.25 per million output tokens on Batch, twice that on the standard API.
% src: common.py OPEN_MODEL_REPO, OPEN_MODEL_REVISION, VLLM_VERSION, LOCAL_MAX_MODEL_LEN, LOCAL_GPU_MEMORY_UTILIZATION, LOCAL_ENGINE_ENV; 04_llm_runner.py run_local_chunk (max-model-len = max(8192, max_completion_tokens + 4096), seed 261, HF_HUB_OFFLINE=1); local_<k>.log (structured_outputs_config backend auto, reasoning_parser openai_gptoss)
\item \textbf{Open-weights arms (Section~\ref{sec:results:open}):} \ident{gpt-oss-20b} from the repository \path{openai/gpt-oss-20b} at revision \path{6cee5e81ee83917806bbde320786a8fb61efebee}, read offline from a local copy, answered by the offline batch runner of vLLM~0.32.0 on one NVIDIA A100-SXM4-80GB GPU per job. The engine uses 22\% of the GPU memory, a maximum model length of 8{,}192 tokens (12{,}096 for a retry at 8{,}000 completion tokens), engine seed 261, its structured-output decoding for the JSON schema with the reasoning parser of gpt-oss, and the PyTorch sampler (\path{VLLM_USE_FLASHINFER_SAMPLER=0}). The request body is that of Appendix~\ref{app:prompt:request} with the model name \texttt{gpt-oss-20b}; no price applies. \ident{llama-3.1-8b-instruct} comes from \path{meta-llama/Llama-3.1-8B-Instruct} at revision \path{0e9e39f249a16976918f6564b8830bc894c89659} in bfloat16, with the same engine settings except 22\% of the GPU memory (27\% for its pilot) and no reasoning parser; its request body omits \ident{reasoning_effort}, and sampling follows the repository's \path{generation_config.json} (temperature 0.6, top-p 0.9).
\end{itemize}

Identical requests need not return identical answers, because no seed is set (and, for the open-weights arms, because the batches that share the engine vary). The test-retest duplicates of Section~\ref{sec:methods:arms} measure this variation, and the deterministic rule, not the model, supplies the reference label of each window.

\subsection{Prompt Text}
\label{app:prompt:body}

The prompt is a single user message. The placeholder \texttt{\{gex\_data\_table\}} stands for the 30 data lines of a window, one per day, in the form \texttt{Day T-29: +3.42B} to \texttt{Day T+0: -12.18B}: a relative day label, the net GEX in billions of US dollars with two decimals and an explicit sign, and neither a date nor a ticker. In the two dated arms (\ident{p34_dated_2024} and \ident{p34_dated_2020}) the label is instead the ISO date of the trading day, as in the November 2025 runs (\texttt{2024-01-02: -33.01B}). The rest of the text is the same for every window, series and symbol (prompt version \ident{v2.0_postgresql}) and is the text of the Phase 4A run.

The listing was checked against the stored requests. With the 30 data lines replaced by the placeholder, it equals the user message of 35{,}365 of 35{,}365 requests (the 35{,}833 requests of Table~\ref{tab:arms} less the 468 requests of the two dated arms). Each stored answer also carries the sha256 of the prompt that produced it in \ident{detection_metadata.prompt_sha256}.
% src: comparison of this listing with the user message of every input_*.jsonl of the non-dated gpt-6-luna runs in data/backups/gex_llm/ (35,365 requests, 35,365 equal; the 40-request synchronous pilot is included; 35,833 ledger requests minus 2 x 234 = 468 requests of the two dated arms = 35,365), 2026-10-31 for the 34,244 requests of the first runs and 2026-11-03 for the 1,121 re-requests (after removing the 30 data lines, one distinct message across all 35,365 requests); also equal to common.build_prompt

\begin{lstlisting}[style=prompt]
You are a market structure analyst specializing in dealer gamma positioning regimes.

TASK: Analyze this 30-day period and determine if it represents a PERSISTENT regime where dealer constraints create forced, directional flows.

## 30-DAY GEX DATA

{gex_data_table}

## REGIME CLASSIFICATION FRAMEWORK

### PERSISTENT REGIMES (Detect These)

**1. PERSISTENT POSITIVE REGIME**
- Definition: Dealers are LONG gamma, forced to sell into strength
- Criteria:
  * >70% of days (21+/30) have positive net GEX
  * Average magnitude >$5B
  * <=5 sign flips across 30 days
  * Stable directional constraint

**Mechanism**: When dealers hold long gamma:
- Price rises -> Dealers MUST sell shares (rebalance)
- Price falls -> Dealers MUST buy shares (rebalance)
- Creates dampening, mean-reverting flows
- Constraint is STRUCTURAL (dealers cannot avoid)

**2. PERSISTENT NEGATIVE REGIME**
- Definition: Dealers are SHORT gamma, forced to buy into strength
- Criteria:
  * >70% of days (21+/30) have negative net GEX
  * Average magnitude >$5B
  * <=5 sign flips across 30 days
  * Stable directional constraint

**Mechanism**: When dealers hold short gamma:
- Price rises -> Dealers MUST buy shares (chase)
- Price falls -> Dealers MUST sell shares (chase)
- Creates amplifying, momentum flows
- Constraint is STRUCTURAL (dealers cannot avoid)

---

### NON-REGIMES (Reject These)

**3. TRANSITIONAL (Reject)**
- Frequent sign flips between positive/negative GEX
- No dominant regime direction (less than 70% same sign)
- Market in regime change period
- Example: 15 positive days, 15 negative days (50/50 split)

**Why Reject**: No persistent constraint. Dealers face mixed conditions daily. Not a structural regime.

**4. LOW CONVICTION (Reject)**
- Consistent sign BUT weak magnitude (<$5B average)
- Example: 25 days positive, avg $2B GEX
- Insufficient constraint to create persistent forced flows

**Why Reject**: Even if sign is consistent, magnitude too weak to force dealers into meaningful positions. Not a structural constraint.

---

## ANALYSIS QUESTIONS

Systematically evaluate the 30-day window:

**Step 1: Sign Persistence**
1. Count days with positive net GEX
2. Count days with negative net GEX
3. Calculate persistence percentage: max(positive_days, negative_days) / 30 * 100
4. Does it meet 70% threshold (21+ days)?

**Step 2: Magnitude Assessment**
1. Calculate average GEX magnitude (absolute value): sum(|net_gex|) / 30
2. Is average magnitude >=$5B?
3. Check for extreme outliers that might distort average

**Step 3: Stability Check**
1. Count sign flips: How many times does GEX switch from pos->neg or neg->pos?
2. Are there <=5 sign flips across 30 days?
3. Stable regime should have low flip count

**Step 4: Regime Classification**
- If Steps 1, 2, 3 all pass AND positive dominates -> PERSISTENT POSITIVE
- If Steps 1, 2, 3 all pass AND negative dominates -> PERSISTENT NEGATIVE
- If Step 1 passes but Step 2 fails -> LOW CONVICTION (reject)
- If Step 1 fails -> TRANSITIONAL (reject)

---

## CONFIDENCE CALIBRATION (Mechanical Guidance)

Use these concrete anchors to calibrate confidence:

**90-100 (Very High Confidence)**
- 25-30 days same sign (83-100% persistence)
- Average magnitude >$10B
- 0-2 sign flips (highly stable)
- Example: "29 negative days, avg $15B, 1 flip"

**70-89 (High Confidence)**
- 21-24 days same sign (70-80% persistence)
- Average magnitude $5-10B
- 2-4 sign flips (moderately stable)
- Example: "23 negative days, avg $7B, 3 flips"

**50-69 (Borderline - Use with Caution)**
- 18-20 days same sign (60-67% persistence)
- Average magnitude $3-5B
- 5-7 sign flips
- Example: "20 negative days, avg $4B, 6 flips"
- **Note**: Borderline cases should generally be REJECTED unless other factors strengthen confidence

**0-49 (Reject - Not Persistent)**
- <18 days same sign (<60% persistence)
- OR average magnitude <$3B
- OR >7 sign flips
- These are NOT persistent regimes

**Important**: Confidence is a FILTER, not a probability. Use it to distinguish clear regimes (70+) from borderline (50-69) from noise (<50).

---

## OUTPUT FORMAT (JSON)

Provide your analysis in this exact JSON structure:

```json
{
    "regime_detected": true/false,
    "regime_type": "persistent_positive|persistent_negative|transitional|low_conviction",
    "positive_days": <count as integer>,
    "negative_days": <count as integer>,
    "avg_magnitude_billions": <value as number>,
    "sign_flips": <count as integer>,
    "persistence_pct": <percentage as number>,
    "confidence": <integer 0-100>,
    "reasoning": "Explain step-by-step why this is/isn't a persistent regime. Reference specific metrics (persistence %, avg magnitude, sign flips). If rejecting, state which criterion failed."
}
```

**IMPORTANT**: All numeric fields (confidence, positive_days, negative_days, sign_flips, avg_magnitude_billions, persistence_pct) MUST be numbers (integers or decimals), NOT words like "thirty-five" or "fifty".

**regime_detected Rules**:
- `true` ONLY if regime_type is "persistent_positive" or "persistent_negative"
- `false` if regime_type is "transitional" or "low_conviction"

---

## KEY PRINCIPLES

1. **Selectivity is Expected**: Most windows will NOT be persistent regimes (expect 30-50% detection rate)

2. **ALL Criteria Must Pass**: Persistence + Magnitude + Stability required for detection

3. **Rejection is Valid**: Saying "no persistent regime" is a correct answer for transitional/weak periods

4. **Mechanical Over Qualitative**: Use concrete thresholds (70%, $5B, 5 flips) rather than subjective judgment

5. **Structural Focus**: Only detect when dealers are FORCED into directional positions by constraints

Analyze the 30-day GEX data above and provide your regime classification in JSON format.
\end{lstlisting}

\subsection{Output Schema}
\label{app:prompt:schema}

The answer is constrained by the following strict JSON schema, passed as \ident{response_format}. With \ident{strict} set, the model can neither omit nor add a field. The boolean \ident{regime_detected} is the primary label of the study and \ident{regime_type} the secondary one; the counts, the mean magnitude and the persistence are the model's own report and are compared with the values of the rule (Section~\ref{sec:results:fidelity}).

\begin{lstlisting}[style=appcode]
{
  "type": "json_schema",
  "json_schema": {
    "name": "regime_classification",
    "strict": true,
    "schema": {
      "type": "object",
      "additionalProperties": false,
      "properties": {
        "regime_detected": {
          "type": "boolean"
        },
        "regime_type": {
          "type": "string",
          "enum": [
            "persistent_positive",
            "persistent_negative",
            "transitional",
            "low_conviction"
          ]
        },
        "positive_days": {
          "type": "integer"
        },
        "negative_days": {
          "type": "integer"
        },
        "avg_magnitude_billions": {
          "type": "number"
        },
        "sign_flips": {
          "type": "integer"
        },
        "persistence_pct": {
          "type": "number"
        },
        "confidence": {
          "type": "integer"
        },
        "reasoning": {
          "type": "string"
        }
      },
      "required": [
        "regime_detected",
        "regime_type",
        "positive_days",
        "negative_days",
        "avg_magnitude_billions",
        "sign_flips",
        "persistence_pct",
        "confidence",
        "reasoning"
      ]
    }
  }
}
\end{lstlisting}

\subsection{Request Body}
\label{app:prompt:request}

Each line of a batch input file is one request of the following form.

\begin{lstlisting}[style=appcode]
{
  "custom_id": "<arm>|<symbol>|<window_end>",
  "method": "POST",
  "url": "/v1/chat/completions",
  "body": {
    "model": "gpt-6-luna",
    "messages": [
      {
        "role": "user",
        "content": "<prompt text of Appendix A.2>"
      }
    ],
    "max_completion_tokens": 4000,
    "reasoning_effort": "medium",
    "response_format": {
      "type": "json_schema",
      "json_schema": "<schema of Appendix A.3>"
    }
  }
}
\end{lstlisting}
